Hlavním cílem tohoto výzkumného úkolu byl pohled do historie umělé inteligence ve hrách a jak se vůbec taková klasická inteligence vytváří. 

Poté byli řečeny příklady 3D akčních her, kde byla umělá inteligence sestrojena trochu jinými zúůsoby, než je tomu v tradiční hře. 

Dále byli popsány čtyři typy strojové učení a u každého byl zmíněn nebo popsán i algoritmus, který pod ten typ strojového učení patří. 

Návrh agenta byla další kapitola, ve které byli zvoleny dva algoritmy, ze kterých by byl agent zkonstruován, aby bylo vyzkoušeno, který z těchto dvou algoritmů je lepší pro AI do her. Byla také navržena vizualizace, kde by agent byl použit. 

Nakonec byli navrženy tři typy hodnocení agenta, kde každé hodnocení ovlivňovat agenta úplně jiným směrem, než ostatní dvě hodnocení.

Sestrojení vizualizace a agenta by proběhlo v diplomové práci, kde by bylo popsáno, jak vůbec byl agent a vizualizace sestrojena a jestli je tato domněnka vůbec sestrojitelná.